\chapter{Wo das Zustandsmodell endet und die Dynamik beginnt}
\label{chap:concept-boundaries}

Teil V hat pose2, pose3, bewegte Frames, Überhöhung, Metric Realization, Soll-
und physischen Zustand sowie Beobachtung bereits eingeführt. Dieses Kapitel
definiert sie nicht erneut. Es grenzt die engeren Verantwortungen von Teil VI
ab, damit die folgenden Kapitel ein einziges Eisenbahnzustandsmodell
konsumieren.

\section{Übernommener Eisenbahnzustand}

Eingabe ist der qualifizierte Zustand aus Teil V an \(s\), einschließlich
Alignment-Bezug, Realisierung, Frame-Konvention, Provenienz, Anwendbarkeit und
Gültigkeit. ``Track State'' ist in den folgenden Kapiteln eine kurze
leserbezogene Bezeichnung für den jeweils anwendbaren Ausschnitt dieser
Struktur, kein konkurrierender Zustandstyp. Ebenso ersetzt ein dargestelltes
Koordinatentupel pose3 nicht, und eine beobachtete Pose ist kein physischer
Zustand.

\section{Neue Verantwortungen}

\begin{description}
\item[Operative Geschwindigkeit] qualifiziert das Durchfahrtsgesetz \(s(t)\)
und seine Ableitungen.
\item[Zeitliche Exposition] komponiert räumliche Krümmungs-, Überhöhungs-,
Profil- und Frame-Größen mit dieser Durchfahrt.
\item[Reduzierte Kinematik] wertet annahmengebundene Indikatoren aus, ohne eine
vollständige Antwort zu behaupten.
\item[Vehicle Space] drückt fahrzeugbezogene Größen relativ zum qualifizierten
Gleis-Frame und einem angegebenen Fahrzeugmodell aus.
\item[Antwortmodellierung] ergänzt die für eine vorhergesagte
Fahrzeug--Gleis-Antwort benötigten physikalischen Parameter.
\item[Dynamisches Gleichgewicht] verwendet eine deklarierte reduzierte
Beziehung unter angegebenen Frame-, Vorzeichen-, Geschwindigkeits- und
Überhöhungskonventionen.
\item[Realisierungsbewusste Dynamik] vergleicht qualifizierte Soll-, realisierte
oder beobachtete Eingaben, ohne sie zusammenzuführen.
\item[Optimierung] variiert deklarierte Freiheitsgrade gegen explizite Ziele
und Nebenbedingungen; sie liefert weder Ziele noch Entscheidungen.
\end{description}

\section{Ein Zustand, mehrere Verwendungen}

Derselbe qualifizierte Eisenbahnzustand kann eine Lichtraumberechnung, einen
reduzierten Querindikator, ein Mehrkörpermodell oder ein
Optimierungsexperiment speisen. Diese Verwendungen erzeugen unterschiedliche
Ausgaben, weil sie unterschiedliche Modelle und Zwecke ergänzen. Sie erzeugen
keine alternativen Definitionen des Alignments oder Eisenbahnzustands.

\begin{principle}[Prinzip der einmaligen Definition]
pose2, pose3, bewegtes Frame, Überhöhung, Realisierung, physischer Zustand und
Beobachtung behalten ihre eingeführten Bedeutungen. Teil VI ergänzt
Durchfahrt, Antwortmodelle und zweckqualifizierte Verwendungen durch Verweis.
\end{principle}

\begin{summary}
Teil VI beginnt bei \(s(t)\). Alles Räumliche wird übernommen; jede zeitliche
oder antwortbezogene Größe muss die zusätzliche Eingabe benennen, die sie
erzeugt.
\end{summary}
